\section{Introduction}
\label{sec:intro}

Counterfactual regret minimization (CFR) \cite{zinkevich2007regret} is the algorithm
of choice for computing equilibria in large games of hidden information, a class that
includes poker, bargaining, and adversarial security settings. Interestingly, tabular
CFR has resisted the move to accelerators, and its fastest implementations still run
on CPUs. Its execution
model explains why. One iteration touches every node of a game tree in millions of tiny
dependent steps. Each step chases a pointer, branches on the node type, and updates a
handful of floats. Every read and write is therefore a data-dependent gather or scatter
\cite{kim2026parallelizing}, and a wide processor finds nothing to batch in such a
walk. Deployed implementations therefore still execute CFR as a node-wise tree
traversal through a generic game interface
\cite{lanctot2019openspiel,liu2024liteefg,steinberger2019pokerrl,li2026real}.

This cost matters because the solves are time-bound. Modern game-playing agents must
solve a game tree with tens of millions of states within a few seconds during play
\cite{burch2014solving,ganzfried2015endgame,moravcik2016refining,brown2017safe,moravvcik2017deepstack,brown2018superhuman,brown2019superhuman,vsustr2019monte,brown2018depth,steinberger2019pokerrl},
so the cost of one solver iteration bounds what any such agent can do. Tabular CFR and
its CFR$^{+}$ refinement \cite{tammelin2014solving} produced the essential solution of
heads-up limit hold'em \cite{bowling2015heads,tammelin2015solving} and drive
superhuman poker AIs
\cite{brown2018superhuman,moravvcik2017deepstack,brown2019superhuman}. The tabular form
remains the field's workhorse: the reference against which every approximation is
checked \cite{lanctot2009monte,brown2019deep,schmid2019variance}, the source of the
exact best-response evaluations that certify solution quality
\cite{johanson2011accelerating,li2026av,li2026agents}, and the inner loop of online re-solving
\cite{burch2014solving,brown2017safe,sustr2020sound,li2025efficient,li2026real}.

That demand has driven many attempts to move the workload onto accelerators, and none has
produced a GPU solver whose iteration time is below that of an efficient CPU
implementation \cite{kim2026parallelizing,DBLP:journals/corr/abs-2508-06559}. The
milestone poker systems sidestepped the GPU entirely and used supercomputer-scale CPU
parallelism instead \cite{bowling2015heads,tammelin2015solving,brown2018superhuman}.
Porting the traversal to a GPU shows why. Every kernel runs for microseconds, so the
time goes into launching kernels and dispatching framework operations
\cite{guide2020cuda,paszke2019pytorch}, and arithmetic is a small fraction of it. The fastest prior GPU CFR
is the sequence-form \cite{koller1997representations} implementation of Kim (2026)
\cite{kim2026parallelizing}. It expresses each tree level as a sparse matrix product,
yet its iteration still takes longer than an optimized CPU solver. Its iteration also
cannot be recorded as a CUDA graph, because the sparse products and per-level index
traffic it relies on are not capturable (\cref{app:profile}). The obstacle lies in the
representation the algorithm runs on, and the workload itself is unusually regular. When the same game is solved for many
iterations, the tree topology, information sets, chance behavior, tensor shapes, and
data dependencies stay constant. Only solver state changes: strategies, reaches,
regrets, values, and the iteration weight. This raises a systems question. How much
of tabular CFR's cost is an artifact of its execution plan, and how much of it can a
compiler remove without touching the algorithm?

\begin{figure}[!htb]
    \centering
    \includegraphics[width=0.62\linewidth]{figures/fig1_hero_scaling.pdf}
    \caption{Steady-state time per CFR iteration versus game size (both axes
    logarithmic; lower is better); backends and variants as in \cref{tab:matrix}.}
    \Description{Log-log scatter plot of steady-state milliseconds per iteration
    against game size in Infosets for six systems. The graph-replayed GPU-CFR curve
    stays nearly flat while baseline curves rise with game size, giving the widest
    margins on the largest games.}
    \label{fig:hero}
\end{figure}

To remove these obstacles, we propose GPU-CFR, which treats a fixed game as a program
to be compiled. GPU-CFR walks a game once (two
players, zero-sum, perfect recall) and emits flat edge,
information-set slot, and depth arrays. An iteration then becomes a fixed sequence of
depth-level batched gather and scatter operations over those arrays. Three techniques
shape the compiled iteration. Static chance folding moves all strategy-independent
chance work to build time. Depth-level execution blocks collapse each traversal stage
into one batched operation per tree level. A dual-lane reach buffer advances both
players' reach probabilities in a single operation, and a sentinel slot removes the
last per-edge branch. The resulting sequence has fixed shapes, indices, and buffer
addresses, so CUDA Graph Replay \cite{guide2020cuda} records it once and replays it
with one graph launch per iteration. The design separates the two sources of run-time
cost. The compiled representation reduces the work issued to the tensor framework, and
graph replay reduces the cost of issuing it.

The compiled iteration pays off most where the trees are largest. \Cref{fig:hero} plots
steady-state time per iteration against game size for every backend we test. The CPU
libraries and the prior GPU implementation slow down as the tree grows. The
graph-replayed GPU-CFR curve stays nearly flat, so the margin widens with scale. On
one A100, a heads-up no-limit Texas hold'em (HUNL) turn subgame with 83{,}040
information sets (Infosets) runs at \turnSteadyMs\,ms per
steady-state CFR$^{+}$ iteration. Against the matched A100 implementation of
Kim (2026)~\cite{kim2026parallelizing}, graph replay is \kimGraphRange$\times$ faster
across the eight-game suite. Against LiteEFG \cite{liu2024liteefg}, one of the
fastest open-source CPU implementations, it is
\liteGraphLargestRange$\times$ faster on the four largest games. Most of this margin
comes from the compiled representation itself: the same compiled solver on eight CPU
threads already beats the A100 baseline by \kimCpuRange$\times$. GPU-CFR in turn
outperforms every CPU path we test, including that eight-thread arm and its one-socket
scaling in \cref{app:profile}.

We answer three systems questions. (Q1) How can a full CFR iteration be compiled into
fixed tensor dataflow while preserving its update semantics? (Q2) Which gains come
from the compiled representation and which come from CUDA Graph Replay? (Q3) When does
the lower iteration cost translate into lower wall-clock time for solution quality and
for repeated solves? Our contributions are:
\begin{enumerate}
    \item \textbf{A game-to-dataflow compiler for tabular CFR.} GPU-CFR maps a fixed
game to flat arrays and a depth-level execution schedule, including static chance
folding and a dual-lane reach representation
(\cref{fig:archpipeline,fig:archdataflow}).
    \item \textbf{CUDA Graph Replay with a causal performance decomposition.} We
separate representation cost from framework dispatch and graph-launch cost through
operation counts, eager/graph ablations, scaling, and repeated-solve measurements.
    \item \textbf{State-of-the-art GPU performance for tabular CFR.} GPU-CFR runs
\kimGraphRange$\times$ faster than the matched A100 baseline across the eight-game
suite and outperforms every tested CPU path (\cref{sec:results}).
\end{enumerate}
